% ============================================================
% Chapter 9: Secure Software Design and Engineering
% Language: Greek — edit freely; regenerate from src/content if preferred.
% ============================================================
\chapter{Ασφαλής Σχεδίαση και Ανάπτυξη Λογισμικού}\label{ch:9}
\chaptermeta{Shift-left και DevSecOps · Έγχυση κώδικα · Cross-site scripting · CSRF και clickjacking · Ασφάλεια API, SSRF και σπασμένος έλεγχος πρόσβασης · Μυστικά και εφοδιαστική αλυσίδα λογισμικού}{Επίπεδο: Μεσαίο επίπεδο · Ανάπτυξη λογισμικού}{Χρόνος μελέτης: 8–10 ώρες}

Οι περισσότερες ευπάθειες δεν δημιουργούνται από τους επιτιθέμενους· γράφονται, ακούσια, από τους προγραμματιστές. Ένας μόνο έλεγχος επικύρωσης που λείπει μπορεί να εκθέσει εκατομμύρια εγγραφές, ενώ η διόρθωση ενός σφάλματος μετά την κυκλοφορία κοστίζει πολύ περισσότερο από την πρόληψή του κατά τον σχεδιασμό. Η ασφαλής μηχανική λογισμικού στοχεύει, επομένως, να καταστήσει την ασφάλεια φυσιολογική ιδιότητα της καθημερινής αναπτυξιακής εργασίας και όχι έναν τελικό έλεγχο πριν από την κυκλοφορία.

Το κεφάλαιο παρουσιάζει πρώτα τη φιλοσοφία \emph{shift-left} και τις αρχές μηχανικής που καθοδηγούν τον ασφαλή σχεδιασμό. Στη συνέχεια εξηγεί, για καθεμιά από τις σημαντικότερες κατηγορίες ευπαθειών ιστού —έγχυση, cross-site scripting, πλαστογράφηση αιτημάτων μεταξύ ιστοτόπων, πλαστογράφηση αιτημάτων από την πλευρά του διακομιστή και σπασμένο έλεγχο πρόσβασης— γιατί προκύπτει το σφάλμα και ποια άμυνα αντιμετωπίζει τη βασική του αιτία. Ολοκληρώνεται με τη διαχείριση μυστικών, την καταγραφή και την προστασία της εφοδιαστικής αλυσίδας λογισμικού. Οι κατηγορίες ευπαθειών που εξετάζονται αντιστοιχούν σε μεγάλο βαθμό στο OWASP Top 10.

\begin{outcomesbox}{Μαθησιακά αποτελέσματα}
Μετά τη μελέτη του κεφαλαίου θα πρέπει να μπορείτε:
\begin{itemize}
  \item Να εξηγείτε την ασφάλεια shift-left και πώς ενσωματώνονται οι δραστηριότητες ασφάλειας σε μια ροή DevSecOps.
  \item Να εξηγείτε τη βασική αιτία της έγχυσης και να εφαρμόζετε παραμετροποιημένα ερωτήματα.
  \item Να διακρίνετε το αποθηκευμένο, το ανακλώμενο και το βασισμένο στο DOM XSS και να εφαρμόζετε κωδικοποίηση εξόδου ανάλογα με το πλαίσιο και CSP.
  \item Να περιγράφετε τις επιθέσεις CSRF, clickjacking, SSRF και τον σπασμένο έλεγχο πρόσβασης, μαζί με τις άμυνές τους.
  \item Να περιγράφετε πώς η ανάλυση σύνθεσης λογισμικού, τα SBOM και το SLSA προστατεύουν την εφοδιαστική αλυσίδα λογισμικού.
\end{itemize}
\end{outcomesbox}

\section{Ασφάλεια shift-left, DevSecOps και βασικές αρχές μηχανικής}\label{sec:9.1}
Ο όρος \textbf{shift-left} σημαίνει τη μετακίνηση των δραστηριοτήτων ασφάλειας νωρίτερα —«προς τα αριστερά»— στο χρονοδιάγραμμα της ανάπτυξης. Οι απαιτήσεις ασφάλειας ορίζονται μαζί με τις λειτουργικές· η \textbf{μοντελοποίηση απειλών} (για παράδειγμα με τη μέθοδο STRIDE: πλαστοπροσωπία, αλλοίωση, αποποίηση, αποκάλυψη πληροφοριών, άρνηση υπηρεσίας, κλιμάκωση προνομίων) εντοπίζει τους κινδύνους όσο ο σχεδιασμός μπορεί ακόμη να αλλάξει με χαμηλό κόστος· και αυτοματοποιημένοι έλεγχοι εκτελούνται σε κάθε αλλαγή του κώδικα. Το \textbf{DevSecOps} ενσωματώνει αυτούς τους ελέγχους στη ροή συνεχούς ενοποίησης και παράδοσης, ώστε μια υποβολή κώδικα που εισάγει μια βιβλιοθήκη με γνωστή ευπάθεια ή έναν ενσωματωμένο κωδικό να αποτυγχάνει στη μεταγλώττιση, όπως ακριβώς και ένας αποτυχημένος έλεγχος μονάδας.

Ορισμένες αρχές μηχανικής καθοδηγούν τη συγγραφή ασφαλούς κώδικα. \textbf{Επικύρωση της εισόδου} στην πλευρά του διακομιστή, με βάση λίστα όσων αναμένονται, επειδή κάθε είσοδος από έξω από το όριο εμπιστοσύνης μπορεί να είναι εχθρική. \textbf{Κωδικοποίηση της εξόδου} σύμφωνα με το πλαίσιο στο οποίο χρησιμοποιείται. \textbf{Ελαχιστοποίηση της επιφάνειας επίθεσης} με την αφαίρεση αχρησιμοποίητων λειτουργιών και σημείων πρόσβασης. \textbf{Ασφαλής αποτυχία}, ώστε τα σφάλματα να απορρίπτουν την πρόσβαση και να μη διαρρέουν εσωτερικές λεπτομέρειες. Και προτίμηση των δοκιμασμένων \textbf{λειτουργιών ασφάλειας των πλαισίων ανάπτυξης} έναντι αυτοσχέδιων μηχανισμών, στο πνεύμα της αρχής της οικονομίας μηχανισμού από το Κεφάλαιο 1.

\section{Έγχυση SQL και NoSQL: παραμετροποίηση έναντι συνένωσης συμβολοσειρών}\label{sec:9.2}
Η \textbf{έγχυση} (injection) συμβαίνει όταν μη αξιόπιστα δεδομένα ερμηνεύονται ως μέρος μιας εντολής. Η κλασική περίπτωση είναι ένα ερώτημα SQL που κατασκευάζεται με συνένωση συμβολοσειρών, όπως \texttt{''SELECT * FROM users WHERE name = ''' + input + '''''}. Αν ένας επιτιθέμενος εισαγάγει \texttt{' OR '1'='1}, η συνθήκη γίνεται πάντα αληθής και το ερώτημα επιστρέφει όλους τους χρήστες· πιο σύνθετα φορτία μπορούν να διαβάσουν άλλους πίνακες, να τροποποιήσουν δεδομένα ή, σε ορισμένες διαμορφώσεις, να εκτελέσουν εντολές του λειτουργικού συστήματος. Η βασική αιτία είναι η \emph{ανάμειξη κώδικα και δεδομένων} σε μία συμβολοσειρά.

Η οριστική άμυνα είναι το \textbf{παραμετροποιημένο ερώτημα} (prepared statement), στο οποίο η δομή του SQL αποστέλλεται στη βάση δεδομένων χωριστά από τις τιμές: \texttt{SELECT * FROM users WHERE name = ?}. Η βάση δεδομένων αντιμετωπίζει τότε την είσοδο αυστηρά ως δεδομένα, όποιους χαρακτήρες κι αν περιέχει. Οι βιβλιοθήκες αντιστοίχισης αντικειμένων–σχέσεων (ORM) παρέχουν την ίδια προστασία όταν χρησιμοποιούνται σωστά. Η επικύρωση εισόδου και οι λογαριασμοί βάσης δεδομένων με ελάχιστα προνόμια είναι πολύτιμα πρόσθετα επίπεδα, όμως η χειροκίνητη διαφυγή ειδικών χαρακτήρων είναι επιρρεπής σε σφάλματα και δεν πρέπει να αποτελεί τη βασική άμυνα. Οι βάσεις NoSQL δεν είναι απρόσβλητες: η αποδοχή αντικειμένων JSON όπως \texttt{\{''\$ne'': null\}} σε ένα ερώτημα MongoDB μπορεί να παρακάμψει την ταυτοποίηση με ακριβώς τον ίδιο τρόπο.

\section{Cross-site scripting: αποθηκευμένο, ανακλώμενο και βασισμένο στο DOM}\label{sec:9.3}
Το \textbf{cross-site scripting (XSS)} είναι έγχυση στον φυλλομετρητή ιστού. Ο επιτιθέμενος κάνει έναν ιστότοπο να παραδώσει κώδικα JavaScript που εκτελείται στον φυλλομετρητή του θύματος με όλα τα προνόμια του ιστοτόπου —διαβάζοντας δεδομένα συνεδρίας, εκτελώντας ενέργειες εκ μέρους του χρήστη ή αλλοιώνοντας τη σελίδα. Στο \textbf{αποθηκευμένο XSS}, το κακόβουλο σενάριο αποθηκεύεται στον διακομιστή (για παράδειγμα σε ένα σχόλιο) και σερβίρεται σε κάθε επισκέπτη. Στο \textbf{ανακλώμενο XSS}, το σενάριο είναι μέρος ενός ειδικά κατασκευασμένου συνδέσμου και επιστρέφεται αμέσως στην απόκριση. Στο \textbf{XSS βασισμένο στο DOM}, η ευπάθεια βρίσκεται εξ ολοκλήρου σε κώδικα της πλευράς του πελάτη που γράφει μη αξιόπιστα δεδομένα στη σελίδα, για παράδειγμα μέσω της ιδιότητας \texttt{innerHTML}.

Η κύρια άμυνα είναι η \textbf{κωδικοποίηση εξόδου ανάλογα με το πλαίσιο}: δεδομένα που εισάγονται σε HTML, σε χαρακτηριστικό HTML, σε JavaScript ή σε URL πρέπει να κωδικοποιούνται σύμφωνα με τους κανόνες του αντίστοιχου πλαισίου. Τα σύγχρονα πλαίσια ανάπτυξης, όπως τα React και Angular, κωδικοποιούν εξ ορισμού, και οι προγραμματιστές πρέπει να αποφεύγουν παρακάμψεις όπως το \texttt{dangerouslySetInnerHTML}, εκτός αν το περιεχόμενο έχει εξυγιανθεί με μια βιβλιοθήκη όπως η DOMPurify. Μια \textbf{Πολιτική Ασφάλειας Περιεχομένου} (Content Security Policy, CSP) προσθέτει άμυνα σε βάθος, δίνοντας εντολή στον φυλλομετρητή να εκτελεί μόνο σενάρια από εγκεκριμένες πηγές, ενώ η σημαία \texttt{HttpOnly} στα cookies εμποδίζει τα σενάρια να διαβάζουν τα cookies συνεδρίας.

\section{Πλαστογράφηση αιτημάτων μεταξύ ιστοτόπων και clickjacking}\label{sec:9.4}
Στην \textbf{πλαστογράφηση αιτημάτων μεταξύ ιστοτόπων} (Cross-Site Request Forgery, CSRF), ένας κακόβουλος ιστότοπος κάνει τον φυλλομετρητή του θύματος να στείλει αίτημα σε άλλον ιστότοπο στον οποίο το θύμα είναι συνδεδεμένο. Επειδή οι φυλλομετρητές επισυνάπτουν αυτόματα τα cookies, ο ιστότοπος-στόχος βλέπει ένα ταυτοποιημένο αίτημα —για παράδειγμα αλλαγής της διεύθυνσης email του λογαριασμού ή μεταφοράς χρημάτων— που ο χρήστης δεν σκόπευε ποτέ να στείλει. Οι άμυνες περιλαμβάνουν τα \textbf{διακριτικά anti-CSRF} (απρόβλεπτες τιμές ενσωματωμένες στις φόρμες, τις οποίες επαληθεύει ο διακομιστής), το χαρακτηριστικό \textbf{SameSite} των cookies, το οποίο εμποδίζει την αποστολή τους με τα περισσότερα αιτήματα μεταξύ ιστοτόπων, και την επανεπαλήθευση ταυτότητας για ευαίσθητες ενέργειες.

Το \textbf{clickjacking} ενσωματώνει αόρατα τον ιστότοπο-στόχο μέσα σε ένα πλαίσιο (frame) στη σελίδα του επιτιθέμενου και εξαπατά τον χρήστη ώστε να πατήσει κουμπιά που δεν βλέπει. Αποτρέπεται με την απαγόρευση της ενσωμάτωσης σε πλαίσια μέσω της οδηγίας \texttt{frame-ancestors} της CSP ή της παλαιότερης κεφαλίδας \texttt{X-Frame-Options}.

\section{Ασφάλεια API: σπασμένος έλεγχος πρόσβασης, SSRF και XXE}\label{sec:9.5}
Ο \textbf{σπασμένος έλεγχος πρόσβασης} είναι η πιο διαδεδομένη κατηγορία του OWASP Top 10. Η συνηθέστερη μορφή του στις διεπαφές API είναι η \emph{μη ασφαλής άμεση αναφορά σε αντικείμενο} (IDOR, γνωστή και ως BOLA): ένα αίτημα όπως \texttt{GET /api/invoices/1043} επιστρέφει το τιμολόγιο ακόμη κι όταν ανήκει σε άλλον πελάτη, επειδή ο διακομιστής ελέγχει \emph{ότι} ο χρήστης είναι συνδεδεμένος αλλά όχι \emph{αν} του ανήκει το αντικείμενο. Κάθε αίτημα πρέπει να εξουσιοδοτείται στον διακομιστή, για κάθε αντικείμενο, σύμφωνα με την αρχή της πλήρους διαμεσολάβησης. Οι διεπαφές API πρέπει επίσης να επιβάλλουν ταυτοποίηση με σωστά επικυρωμένα διακριτικά, να εφαρμόζουν όρια ρυθμού αιτημάτων και να επιστρέφουν μόνο τα πεδία που χρειάζεται ο πελάτης.

Η \textbf{πλαστογράφηση αιτημάτων από την πλευρά του διακομιστή} (Server-Side Request Forgery, SSRF) εξαπατά έναν διακομιστή ώστε να πραγματοποιήσει αιτήματα για λογαριασμό του επιτιθέμενου —για παράδειγμα, μια λειτουργία «λήψη εικόνας από URL» που κατευθύνεται σε εσωτερικές υπηρεσίες ή στο σημείο μεταδεδομένων του νέφους \texttt{169.254.169.254}, το οποίο μπορεί να αποκαλύψει προσωρινά διαπιστευτήρια. Οι άμυνες περιλαμβάνουν λίστες επιτρεπόμενων προορισμών, αποκλεισμό εσωτερικών περιοχών διευθύνσεων και χρήση της θωρακισμένης υπηρεσίας μεταδεδομένων (IMDSv2). Οι επιθέσεις \textbf{εξωτερικής οντότητας XML} (XML External Entity, XXE) εκμεταλλεύονται αναλυτές XML που επιλύουν εξωτερικές αναφορές, επιτρέποντας την ανάγνωση αρχείων από τον διακομιστή· η θεραπεία είναι η απενεργοποίηση των εξωτερικών οντοτήτων και των ορισμών τύπου εγγράφου στη ρύθμιση του αναλυτή.

\section{Μυστικά, καταγραφή και εφοδιαστική αλυσίδα λογισμικού}\label{sec:9.6}
Τα \textbf{μυστικά} —κωδικοί, κλειδιά API, ιδιωτικά κλειδιά— δεν πρέπει ποτέ να γράφονται σε πηγαίο κώδικα ή σε αρχεία ρυθμίσεων που καταχωρίζονται στο σύστημα ελέγχου εκδόσεων, επειδή τα αποθετήρια αντιγράφονται, κοινοποιούνται και ενίοτε δημοσιοποιούνται. Φυλάσσονται σε ειδικό διαχειριστή μυστικών (όπως το HashiCorp Vault ή ένα θησαυροφυλάκιο κλειδιών στο νέφος), εισάγονται κατά την εκτέλεση και εναλλάσσονται τακτικά· εργαλεία σάρωσης μυστικών στη ροή ανάπτυξης εντοπίζουν τις ακούσιες καταχωρίσεις. Η \textbf{καταγραφή ασφάλειας} πρέπει να καταγράφει συμβάντα ταυτοποίησης, αποτυχίες ελέγχου πρόσβασης και διαχειριστικές ενέργειες με επαρκές πλαίσιο για τη διερεύνηση, αλλά δεν πρέπει ποτέ να καταγράφει κωδικούς, διακριτικά ή περιττά προσωπικά δεδομένα.

Οι σύγχρονες εφαρμογές αποτελούνται κυρίως από κώδικα τρίτων: ένα τυπικό έργο μπορεί να εξαρτάται από εκατοντάδες πακέτα ανοιχτού κώδικα. Η \textbf{ανάλυση σύνθεσης λογισμικού} (Software Composition Analysis, SCA) εντοπίζει αυτές τις εξαρτήσεις και επισημαίνει όσες έχουν γνωστές ευπάθειες, όπως στο περιστατικό Log4Shell του 2021. Ένας \textbf{κατάλογος υλικών λογισμικού} (Software Bill of Materials, SBOM), σε μορφές όπως SPDX ή CycloneDX, απαριθμεί κάθε στοιχείο, ώστε τα επηρεαζόμενα συστήματα να εντοπίζονται μέσα σε λίγα λεπτά όταν ανακοινώνεται μια νέα ευπάθεια. Το πλαίσιο \textbf{SLSA} (\emph{Supply-chain Levels for Software Artifacts}) ορίζει αυξανόμενα επίπεδα διασφάλισης για τη διαδικασία κατασκευής λογισμικού —όπως υπογεγραμμένες, αναπαραγώγιμες εκδόσεις με επαληθεύσιμη προέλευση— ώστε επιθέσεις όπως η παραβίαση της διαδικασίας κατασκευής της SolarWinds να γίνονται ανιχνεύσιμες.

\section*{Βασικοί όροι}
\begin{description}[style=nextline,leftmargin=1.2em]
  \item[Shift-left] Μετακίνηση των δραστηριοτήτων ασφάλειας σε πρωιμότερα στάδια του κύκλου ανάπτυξης.
  \item[Μοντελοποίηση απειλών] Συστηματικός εντοπισμός απειλών κατά ενός σχεδιασμού, π.χ. με τη μέθοδο STRIDE.
  \item[Παραμετροποιημένο ερώτημα] Ερώτημα με σταθερή δομή, του οποίου οι τιμές διαβιβάζονται χωριστά ως δεδομένα.
  \item[XSS] Cross-site scripting: έγχυση σεναρίου που εκτελείται στον φυλλομετρητή του θύματος υπό την προέλευση του ιστοτόπου.
  \item[IDOR / BOLA] Πρόσβαση σε αντικείμενα μέσω αναγνωριστικού χωρίς έλεγχο αν ο αιτών είναι εξουσιοδοτημένος γι' αυτά.
  \item[SBOM] Κατάλογος υλικών λογισμικού: απογραφή όλων των στοιχείων ενός λογισμικού.
\end{description}

\section*{Σύνοψη κεφαλαίου}
\begin{itemize}
  \item Η ασφάλεια κοστίζει λιγότερο όταν ενσωματώνεται νωρίς: η μοντελοποίηση απειλών και οι αυτοματοποιημένοι έλεγχοι στη ροή ανάπτυξης αποτελούν τον πυρήνα του DevSecOps.
  \item Η έγχυση προκύπτει από την ανάμειξη κώδικα και δεδομένων· τα παραμετροποιημένα ερωτήματα εξαλείφουν τη βασική αιτία.
  \item Το XSS αποτρέπεται με κωδικοποίηση εξόδου ανάλογα με το πλαίσιο, ασφαλείς προεπιλογές των πλαισίων ανάπτυξης και CSP.
  \item Τα CSRF, clickjacking, SSRF και XXE έχουν συγκεκριμένες, καθιερωμένες άμυνες· ο σπασμένος έλεγχος πρόσβασης απαιτεί εξουσιοδότηση κάθε αντικειμένου στην πλευρά του διακομιστή.
  \item Οι διαχειριστές μυστικών, η προσεκτική καταγραφή, η SCA, τα SBOM και το SLSA προστατεύουν την εφαρμογή και την εφοδιαστική της αλυσίδα.
\end{itemize}

\section*{Ερωτήσεις ανασκόπησης}
\begin{enumerate}
  \item Εξηγήστε γιατί η διαφυγή των εισαγωγικών είναι ασθενέστερη άμυνα κατά της έγχυσης SQL σε σύγκριση με τα παραμετροποιημένα ερωτήματα.
  \item Κατατάξτε ως αποθηκευμένο, ανακλώμενο ή βασισμένο στο DOM XSS την εξής περίπτωση: μια σελίδα αναζήτησης που εμφανίζει το ερώτημα από το URL χρησιμοποιώντας innerHTML.
  \item Γιατί το χαρακτηριστικό SameSite των cookies περιορίζει το CSRF αλλά όχι το XSS;
  \item Όταν ανακοινώθηκε το Log4Shell, πώς θα βοηθούσε ένα SBOM έναν οργανισμό να ανταποκριθεί;
\end{enumerate}
